\section{Use a stopped boundary as evidence}\label{transpose-p-boundaries}
A breakpoint at a source label stops before the instruction at that label. At
\texttt{transpose\_done}, the return value and output can therefore be checked
before the caller changes temporary registers or stores metadata. This is also
the place to check the saved registers and restored stack pointer. Checking only
after program exit can miss a caller that accidentally hides a callee's error.

The automated checks install their data before starting the observation. They
resolve data addresses without executing simulated reads, run the selected kernel,
and stop at the return boundary. A complete access log then distinguishes a
matrix read from a debugger inspection and a word access from a byte access.
Protected-region store traces detect forbidden writes even if the program later
restores the old value. Finally, the checks continue to exit and inspect memory
again, so code after the observed return cannot silently corrupt the result.

These boundaries make the observations repeatable without fixing an internal
instruction count or register assignment. The book's particular scratch-register
states explain its walkthrough. The submission contract specifies which states
must be externally observable for any equivalent implementation.
